\begin{theorem}[The actual contragredient regular charge]
    \label{thm:peps_regular_contragredient_charge}
    \lean{TNLean.PEPS.regularChargeLabels_inv_mem}
    \leanok
    \uses{thm:peps_regular_charge_labels, thm:peps_irreducible_dual}
    For every actual regular charge character
    $\chi$, the function $\chi^\vee(g)=\chi(g^{-1})$ is also an actual
    regular charge character. For unitary characters this is $\overline\chi$.
    This is the contragredient charge in
    \cite[charge-pair creation]{Schuch2010PEPS}, lines 2514--2518.
\end{theorem}
\begin{proof}\leanok
    The contragredient representation is irreducible, and its character
    is $g\mapsto\chi(g^{-1})$. Every irreducible character occurs in
    the actual regular representation.
\end{proof}

\begin{definition}[The contragredient label of an actual regular charge]
    \label{def:peps_regular_dual_charge_label}
    \lean{TNLean.PEPS.regularDualChargeLabel}
    \leanok
    \uses{thm:peps_regular_contragredient_charge}
    Label $\chi^\vee$ by its function $g\mapsto\chi(g^{-1})$ in the actual
    finite family of regular irreducible characters.
\end{definition}

\begin{theorem}[The two original-spin charge slices of the correlated pair]
    \label{thm:peps_regular_physical_charge_pair_slices}
    \lean{TNLean.PEPS.regularChargePairCoefficient_top,
      TNLean.PEPS.regularChargePairCoefficient_middle,
      TNLean.PEPS.exists_regularPhysicalChargePairSliceMeasurement}
    \leanok
    \uses{def:peps_regular_dual_charge_label,
      thm:peps_regular_contragredient_charge}
    In the correlated coefficient $\chi(pm^{-1}t)$, fixing $m$ gives the
    translated $\chi$-character on $t$. Fixing $t$ gives
    \[
        \chi(pm^{-1}t)=\chi^\vee((tp)^{-1}m).
    \]
    For two original regular $G$-isometric site maps sharing one bond,
    one complete positive orthogonal projective measurement distinguishes
    both types of character-weighted contraction. It is chosen before
    all charges, $p$, fixed partner labels and remaining half-edge labels.
    The first slice has outcome $\chi$, and the second has outcome
    $\chi^\vee$; every other outcome annihilates the corresponding column.

    This is an auxiliary shared-bond slice statement. The regional assertion is derived separately by literal contraction
    slicing along the two internal bonds; no regional factorization is
    assumed in this slice statement.
\end{theorem}
\begin{proof}\leanok
    Character cyclicity identifies the second slice with the translated
    dual character. Apply the already derived complete original-spin
    charge measurement to each literal shared-bond column.
\end{proof}
